\documentclass[12pt]{report}
\usepackage[a4paper, margin=2.5cm]{geometry}
\pagenumbering{gobble}
\usepackage{hyperref}
\usepackage{tcolorbox}
\newcommand{\student}[3]{\begin{tcolorbox}\noindent Introduction to
    Psycholinguistics (WS2021-2022) \\ Report of \textbf{Paper~#3} by
    \textbf{#1 (#2)} 
  \end{tcolorbox}}



\begin{document}

%% Change:
\student{Shawon Ashraf}{3460184}{5}


\paragraph{APA reference:} 
McRae, K., Hare, M., Elman, J. L., & Ferretti, T. (2005). A basis for generating expectancies for verbs from nouns. \textit{Memory & cognition}, 33(7), 1174-1184.


\paragraph{Research Question(s):}
\begin{enumerate}
    \item Do elements other than verbs play a role in generating expectancies?
    \item Is event-based knowledge connected to expectancy generation?
\end{enumerate}
\vfill
\paragraph{Methodology:} 
The authors use two experiments in their work, both based on the notion of testing whether a single noun can lead to sufficient information required for event-knowledge. The experiments are also based word-word priming only.


\noindent
\\ For the first experiment, short stimulus onset asynchrony (SOA) is used to gain insights about the organization of semantic memory. There were forty participants in the experiment and all of them were native speakers of English. The participants were provided two lists with equally divided number of related and unrelated primes to gauge their ability to comprehend events based on a theme. Each participant was subjected to silent reading words on a computer screen and speak out every second word quickly and accurately. Their response time were recorded in the process, alongside pronunciation errors. 

\noindent
\\ The second experiment has the same task setup as the first one. However, in this case, long SOA is used instead of short SOA. There were sixty eight participants in the experiment, again, all of them being native speakers of English.


\vfill

\paragraph{Results:} 
The results of the first experiment indicate that prior knowledge of an event allows a comprehender to generate expectancies accordingly for yet to be seen concepts in a language stream. Also, the results were based on noun-verb primed pairs, which answer the first research question. The second experiment had similar results as the first one. The goal of this experiment was to test whether verbs actually dictate event to expectancy generation or whether it is that prior knowledge of verbs as role fillers take precedence. 

\vfill

\paragraph{Personal Opinion:} 
What I found interesting about this paper is that it addresses the concept that in an event driven situation, verbs will take precedence towards understanding and comprehension of language. However, authors omit multi word priming, which would have been an interesting case study on the presented idea. Also, prior comprehension knowledge of a language user can be a factor here, which was not checked in the paper. Different L1 users of the same language can comprehend concepts differently. In my view this could have influenced the results in the experiments.
\vfill


\end{document}

